\documentclass{article}
\usepackage{amsmath,amssymb}
\usepackage{algorithm}
\usepackage{algorithmic}

\begin{document}

\section{Markov Decision Processes}
\quad A Markov Decision Process (MDP) is a mathematical framework used to model decision-making in situations where outcomes are partly random and partly under the control of a decision maker. An MDP is
defined by the tuple \( (S, A, P, R, \gamma) \), where:
\begin{itemize}
    \item \( S \) is a finite set of states.
    \item \( A \) is a finite set of actions.
    \item \( P(s' | s, a) \) is the state transition probability function, which gives the probability of transitioning to state \( s' \) from state \( s \) after taking action \( a \).
    \item \( R(s, a) \) is the reward function, which gives the immediate reward received after taking action \( a \) in state \( s \).
    \item \( \gamma \in [0, 1) \) is the discount factor, which determines the importance of future rewards.
\end{itemize}

A policy \( \pi(a | s) \) defines the probability of taking action \( a \) when in state \( s \).

The return \( G_t \) is defined as the total discounted reward from time step \( t \) onward:
\begin{equation}
G_t = \sum_{k=0}^{\infty} \gamma^k R_{t+k+1}
\end{equation}

The goal of solving an MDP is to find a policy \( \pi: S \rightarrow A \) that maximizes the expected cumulative return from each state.

Note: All of the above assume a discrete state and action space, as well as discrete time steps.

\section{Bellman Equations}

The value function \( V_\pi(s) \) for a given policy \( \pi \) is defined as the expected return when starting in state \( s \) and following policy \( \pi \):

\begin{equation}
V_\pi(s) = \mathbb{E}_\pi [G_t | S_t = s]
\end{equation}

The action-value function \( Q_\pi(s, a) \) is defined as the expected return when starting in state \( s \), taking action \( a \), and thereafter following policy \( \pi \):

\begin{equation}
Q_\pi(s, a) = \mathbb{E}_\pi [G_t | S_t = s, A_t = a]
\end{equation}

The action-value function is often referred to as the Q-function.


Bellman equations provide a recursive relationship for the value functions. The Bellman equation for the value function \( V_\pi(s) \) is given by:

\begin{equation}
V_\pi(s) = \sum_a \pi(a | s) \sum_{s', r} P(s', r | s, a) [r + \gamma V_\pi(s')]
\end{equation} 

The derivation of the Bellman equation for the value function is as follows:

\begin{align*}
V_\pi(s) &= \mathbb{E}_\pi [G_t | S_t = s] \\
&= \mathbb{E}_\pi [R_{t+1} + \gamma G_{t+1} | S_t = s] \\ 
&= \sum_a \pi(a | s) \sum_{s', r} P(s', r | s, a) [r + \gamma V_\pi(s')]
\end{align*}

\section{Policy Iteration}
\label{sec:policy_iteration}

Policy iteration is an algorithm used to find the optimal policy for a Markov Decision Process (MDP). It consists of two main steps: policy evaluation and policy improvement. The algorithm iteratively evaluates the current policy and then improves it until convergence to the optimal policy.

\begin{algorithm}
\caption{Policy Iteration}
\begin{algorithmic}[1]
\STATE Initialize policy \( \pi \) arbitrarily
\REPEAT
    \STATE \textbf{Policy Evaluation:} Compute \( V_\pi(s) \) for all states \( s \) using the Bellman equation
    \begin{equation}
    V_\pi(s) = \sum_a \pi(a | s) \sum_{s', r} P(s', r | s, a) [r + \gamma V_\pi(s')]
    \end{equation}
    \STATE \textbf{Policy Improvement:} Update policy \( \pi \) using the greedy policy derived from the value function
    \begin{equation}
    \pi'(a | s) = \begin{cases}
    1 & \text{if } a = \arg\max_a Q_\pi(s, a) \\
    0 & \text{otherwise}
    \end{cases}
    \end{equation}
\UNTIL{Convergence}
\end{algorithmic}
\end{algorithm}

\end{document}